\documentclass[12pt, a4paper]{article}
\usepackage[margin=2cm]{geometry}
\usepackage{amsmath}
\usepackage{amssymb}
\usepackage{graphicx} % For images, not used directly for this content but required
\usepackage{float}    % For image/table placement, not used directly but required
\usepackage{booktabs} % For professional tables, not used directly but required
\usepackage{array}    % For extended array capabilities, not used directly but required
\usepackage{tabularx} % For tables with fixed width, not used directly but required
\usepackage{longtable}% For multi-page tables, not used directly but required
\usepackage{multirow} % For multirow cells in tables, not used directly but required
\usepackage{ctex}     % Must only use ctex for Chinese support, no fontset parameter

\begin{document}

\section*{1. 投影矩阵的性质}

对 $P \in \mathbb{R}^{n \times n}$，若它满足 $P^2=P$，称其为投影矩阵，证明：
\begin{enumerate}
    \item $I-P$ 为投影矩阵
    \item $R(I-P) = N(P)$, $N(I-P) = R(P)$
    \item $R(P) \cap N(P) = \{0\}$
\end{enumerate}

\subsection*{证明:}

\begin{enumerate}
    \item \textbf{证明 $I-P$ 为投影矩阵}
    由于 $P$ 为投影矩阵，则有 $P^2 = P$。
    考虑 $(I-P)^2$:
    \begin{align*}
    (I-P)^2 &= I^2 - IP - PI + P^2 \\
             &= I - P - P + P \\
             &= I - P
    \end{align*}
    因为 $(I-P)^2 = I-P$，所以 $I-P$ 也是一个投影矩阵。

    \item \textbf{证明 $R(I-P) = N(P)$ 和 $N(I-P) = R(P)$}
    \begin{itemize}
        \item \textbf{证明 $R(I-P) = N(P)$:}
        \begin{itemize}
            \item \textbf{证明 $R(I-P) \subseteq N(P)$:}
            设 $y \in R(I-P)$。则存在 $x \in \mathbb{R}^n$ 使得 $y=(I-P)x$。
            对 $y$ 左乘 $P$ 得到：
            $Py = P(I-P)x = (P-P^2)x = (P-P)x = 0x = 0$。
            故 $y \in N(P)$。因此，$R(I-P) \subseteq N(P)$。

            \item \textbf{证明 $N(P) \subseteq R(I-P)$:}
            设 $x \in N(P)$。则 $Px = 0$。
            我们可以将 $x$ 写成 $x = IX = (I-P)x + Px$。
            由于 $Px=0$，所以 $x = (I-P)x + 0 = (I-P)x$。
            故 $x \in R(I-P)$。因此，$N(P) \subseteq R(I-P)$。
        \end{itemize}
        综合以上两点，得出 $R(I-P) = N(P)$。

        \item \textbf{证明 $N(I-P) = R(P)$:}
        \begin{itemize}
            \item \textbf{证明 $R(P) \subseteq N(I-P)$:}
            设 $y \in R(P)$。则存在 $x \in \mathbb{R}^n$ 使得 $y=Px$。
            对 $y$ 左乘 $(I-P)$ 得到：
            $(I-P)y = (I-P)Px = (P-P^2)x = (P-P)x = 0x = 0$。
            故 $y \in N(I-P)$。因此，$R(P) \subseteq N(I-P)$。

            \item \textbf{证明 $N(I-P) \subseteq R(P)$:}
            设 $x \in N(I-P)$。则 $(I-P)x = 0$。
            这表示 $x - Px = 0$，即 $x = Px$。
            故 $x \in R(P)$ (因为 $x$ 可以表示为 $P$ 乘以某个向量 $x$ 本身)。
            因此，$N(I-P) \subseteq R(P)$。
        \end{itemize}
        综合以上两点，得出 $N(I-P) = R(P)$。
    \end{itemize}

    \item \textbf{证明 $R(P) \cap N(P) = \{0\}$}
    设 $x \in R(P) \cap N(P)$。
    由 $x \in N(P)$，可知 $Px = 0$。
    由 $x \in R(P)$，可知存在 $y \in \mathbb{R}^n$ 使得 $x = Py$。
    将 $x=Py$ 代入 $Px=0$ 中，得到 $P(Py)=0$，即 $P^2y=0$。
    由于 $P$ 是投影矩阵，$P^2=P$，所以 $Py=0$。
    因为 $x=Py$，所以 $x=0$。
    因此，$R(P) \cap N(P) = \{0\}$。
\end{enumerate}

\section*{2. 最小二乘问题}

计算最小二乘问题 $\min \|Ax-b\|^2$，其中
\[ A=\begin{pmatrix}
1 & -1 & 1 \\
1 & 3 & 3 \\
1 & 3 & 5 \\
1 & 3 & 7
\end{pmatrix}, \quad b=\begin{pmatrix}
1 \\
1 \\
1 \\
1
\end{pmatrix} \]

\subsection*{解:}
首先，通过验证矩阵 $A$ 的列向量是线性无关的，可知 $R(A)=3$。
最小二乘解 $\hat{x}$ 满足正规方程 $A^TAx = A^Tb$。

首先计算 $A^TA$:
\begin{align*}
A^TA &= \begin{pmatrix}
1 & 1 & 1 & 1 \\
-1 & 3 & 3 & 3 \\
1 & 3 & 5 & 7
\end{pmatrix}
\begin{pmatrix}
1 & -1 & 1 \\
1 & 3 & 3 \\
1 & 3 & 5 \\
1 & 3 & 7
\end{pmatrix} \\
&= \begin{pmatrix} % Transcribing directly from image's calculation result
4 & 8 & 4 \\
8 & 20 & 24 \\
4 & 24 & 84
\end{pmatrix}
\end{align*}

接下来计算 $A^Tb$:
\begin{align*}
A^Tb &= \begin{pmatrix}
1 & 1 & 1 & 1 \\
-1 & 3 & 3 & 3 \\
1 & 3 & 5 & 7
\end{pmatrix}
\begin{pmatrix}
1 \\
1 \\
1 \\
1
\end{pmatrix} \\
&= \begin{pmatrix}
1 \cdot 1 + 1 \cdot 1 + 1 \cdot 1 + 1 \cdot 1 \\
-1 \cdot 1 + 3 \cdot 1 + 3 \cdot 1 + 3 \cdot 1 \\
1 \cdot 1 + 3 \cdot 1 + 5 \cdot 1 + 7 \cdot 1
\end{pmatrix} \\
&= \begin{pmatrix}
4 \\
8 \\
16
\end{pmatrix}
\end{align*}

正规方程 $A^TAx = A^Tb$ 变为：
\[ \begin{pmatrix}
4 & 8 & 4 \\
8 & 20 & 24 \\
4 & 24 & 84
\end{pmatrix}
\begin{pmatrix}
x_1 \\
x_2 \\
x_3
\end{pmatrix} =
\begin{pmatrix}
4 \\
8 \\
16
\end{pmatrix} \]
解此线性方程组，得到：
\[ x = \begin{pmatrix}
-2 \\
1 \\
0
\end{pmatrix} \]

\end{document}